\statement{B.81}{Proved}{A smaller genuine bundle and whole filling kernels inside specified free factors}
The actual marginal comparison of B.78 admits a genuine seven-puncture
bundle over a six-punctured plane, with source group \(F_{10}\) filling
onto \(F_6\). Its complete based monodromy and four source meridians are
specified below. The source kernel image contains the entire filling
kernels of two specified rank-five free factors and two specified
rank-two free factors of the actual fiber \(F_7\). These are exact
intersection statements, also for actual compact images in full
completions. They do not assert normality of the whole source image.

\begin{proof}
\emph{The actual unordered quotient and its two covers.}\par
Use the normalized ordered surface of B.78 and put
\[
 \begin{gathered}
 q=vw,\qquad T=\frac{v+w}{2},\\
 X=\{(q,T)\in\mathbb C^2:q\notin\{0\}\cup\mu _5,
       \ T^2\ne q,\ q-2uT+u^2\ne0\ (u\in\mu _5)\}.
 \end{gathered}
 \tag{B.81.1}
\]
The swap of \(v,w\) acts freely, since the diagonal is deleted.
Its quotient is exactly \(X\), with inverse unordered roots
\(T\pm\sqrt{T^2-q}\). The five lines in (B.81.1) are tangent to
\(q=T^2\) at \((T,q)=(u,u^2)\). The four lines with \(u\ne1\)
are precisely the deleted set \(\mathcal U\); the line with \(u=1\)
is retained.

The original positive marginal has the other actual double cover
\[
       s^2=q,\qquad T=s\frac{1+p}{1-p}.
 \tag{B.81.2}
\]
Its seven fiber punctures become
\(T=\pm s,(s^2+u^2)/(2u)\). The positive source
\(p=(1-s)/(1+s)\) becomes \(T=1\). Both covering characters
extend after filling \(\mathcal U\): they are the mod-two winding
characters of \(q\) and \(T^2-q\), whose zero divisors remain
retained. On the source they are the characters of \(q\) and
\(1-q\). Thus the ambient filling kernel and the source filling
kernel both lie in the two corresponding covering subgroups.
The covering embeddings identify their actual images. Finite-index
subgroups carry their full induced profinite topology, so the same
kernel comparisons hold after full completion. This uses actual
closed ambient kernels, not abstract completions of infinitely
generated discrete kernels. It is the same finite-cover comparison
as B.78, now with the further free unordered quotient made explicit.

\emph{A genuine bundle, including the actual source.}\par
Projection onto \(q\) has base
\(B=\mathbb C\setminus(\{0\}\cup\mu _5)\) and seven fiber
punctures
\[
       \sqrt q,\ -\sqrt q,\qquad
       T_u(q)=\frac{q+u^2}{2u}\quad(u\in\mu _5).
 \tag{B.81.3}
\]
The square roots meet only at zero. For distinct \(u,v\), the
line points meet exactly at \(q=uv\in\mu _5\), while
\[
                    T_u(q)^2-q=\frac{(q-u^2)^2}{4u^2}.
 \tag{B.81.4}
\]
Consequently no collision occurs over \(B\). Local analytic root
choices and disjoint-disk isotopies give bundle trivializations.
The continuous far section \(T=|q|+2\) avoids all punctures.
The group is therefore \(F_7\rtimes F_6\), with thirteen generators
and forty-two action relations after the specified six base lifts
are chosen. These counts are not minimality claims.

After filling the four \(\mathcal U\)-points, the actual source
\(T=1\) is a global section: its possible collisions with
\(\pm\sqrt q\) and \((q+1)/2\) occur only at the already removed
value \(q=1\). Before filling, its four additional holes are
\[
               q_u=2u-u^2\qquad(u\in\mu _5\setminus\{1\}).
 \tag{B.81.5}
\]
They are distinct, since their pair difference is
\((u-v)(2-u-v)\), and \(|q_u|=|2-u|>1\), so none is an old
base hole. Thus the source is \(F_{10}\), filling onto \(F_6\).
Its two source double covers have ranks nineteen, with filled ranks
eleven, as required by the preceding kernel comparison.

Choose \(q_*=1/100,T_*=1\). Both roots have modulus \(1/10\),
and every line point has modulus at most \(101/200\), so the
source point is a literal exterior basepoint. The source transport
from B.78's already specified normalized endpoint is
\[
 -9999\longrightarrow-9999+100i\longrightarrow100i
                       \longrightarrow1/100,\qquad T=1.
 \tag{B.81.6}
\]
The exact distance audit verifies avoidance of all ten source holes.
The root covers lift this actual path; no unspecified change of the
source subgroup is introduced.

Write \(u_j=\xi^j\), where \(\xi=e^{2\pi i/5}\). Label the five
line points by \(0,\ldots,4\), the source by five, and the two roots
by six and seven. For the projection \((1+371i/1000)T\), the initial
order is
\((2,3,7,1,6,4,0,5)\). The source is last. Let
\(x_1,\ldots,x_7\) be the resulting based fiber meridians.
The retained indices are \(3,5,7\), and the deleted indices are
\(1,2,4,6\). The four inverse source-meridian images are the
following literal words:
\[
\begin{split}
 w_4&=x_6^{-1}x_4x_6,\\
 w_1&=x_6^{-1}x_3^{-1}x_2^{-1}x_1x_2x_3x_6,\\
 w_2&=x_7^{-1}x_6^{-1}x_5^{-1}x_2x_5x_6x_7,\\
 w_6&=x_7^{-1}x_6^{-1}x_5^{-1}x_6x_5x_6x_7.
\end{split}
 \tag{B.81.7}
\]
Let \(\alpha_0\) denote the based action around zero, and
\(\alpha_{j+1}\) the action around \(u_j^2\), \(0\leq j<5\).
The full six automorphisms, their inverses and the four words are
in \path{work_DE/next_marginal_source/tangent_bundle_fiber_data.json}.
This finite table specifies the presentation without relying on local
conjugacy types. The local types themselves are a half twist at zero,
and, at each \(u_j^2\), a fourth power on the tangent pair and two
disjoint full twists on line pairs. The latter have exponent eight
on the seven-point fiber. At \(q=1\) the stationary source also
participates in the source-based eight-strand motion; that distinction
is retained in the full calculation.

The producers \path{tangent_bundle_braid.py} and
\path{extract_tangent_bundle_fiber.py}, in the same directory, retain
their genuine outputs and all fifty segment certificates for the ten
distinguished lassos. The audit \path{audit_tangent_bundle.py}
checks completeness in \(\mathbb Q[\xi]\), actual coefficient
deformation, disjoint critical squares, the common source path and
the distinguished whiskers. The last-strand conjugator is
antihomomorphic; its inverse is the fiber homomorphism used in
(B.81.7). All six actions and inverse actions are checked on a whole
free fiber basis by literal Artin conjugation. The independent check
\path{work_CE/code/verify_tangent_bundle_all_actions.py}, with its
actual output, verifies all eighty-four such comparisons using a
separate free-word evaluator.

\emph{The actual source image and a source-contained free basis.}\par
Let
\[
 R=\ker(F_7\longrightarrow F_3),\qquad
 I=\operatorname{im}\bigl(\ker(F_{10}\longrightarrow F_6)
                                      \longrightarrow F_7\bigr).
 \tag{B.81.8}
\]
The six old source meridians split the source filling, since the
retained group is free. Thus \(I\) is exactly the subgroup generated
by all six-action translates of the four words (B.81.7). In particular
it is invariant under every \(\alpha_i^{\pm1}\), and \(I\leq R\).

Make the following actual free-basis changes:
\[
\begin{gathered}
 U_1=x_3^{-1}x_2^{-1}x_1x_2x_3,\quad
 U_2=x_2,\quad U_4=x_4,\quad U_6=x_6,\\
 A=U_6^{-1}x_3U_6,\quad B=U_6^{-1}x_5U_6,\quad C=x_7,\qquad
 U=\langle U_1,U_2,U_4,U_6\rangle,\quad D=BC.
\end{gathered}
 \tag{B.81.9}
\]
The inverse substitutions are checked literally. The finite certificate
\path{tangent_nielsen_marking.py} proves \(U\leq I\), with explicit
source words: it reevaluates the fifty-two first-stage orbit words and
checks every returned free-group expansion. Moreover all four original
words belong to \(\langle U,\alpha_1(U)\rangle\). Hence \(I\) is
also exactly the subgroup generated by the six-action translates of
this entire free factor \(U\).

\emph{A whole rank-five kernel.}\par
The exact action table gives
\[
 \alpha_1(D)=D,\qquad
 \alpha_1(U)=D^{-1}UD.
 \tag{B.81.10}
\]
The internal automorphism after removing the displayed conjugator
conjugates \(\langle U_1,U_2\rangle\) by \(U_2U_1\), and
\(\langle U_4,U_6\rangle\) by \(U_4U_6\). It is an actual
free-factor automorphism. Consequently
\(D^nUD^{-n}\leq I\) for every integer \(n\).
Since \(D=BC\) is primitive in the retained free factor,
\(V_D=U*\langle D\rangle\) is a rank-five free factor of \(F_7\).
The upper bound \(I\leq R\) now gives the exact equality
\[
 I\cap V_D
    =\ker\bigl(U*\langle D\rangle\longrightarrow\langle D\rangle\bigr).
 \tag{B.81.11}
\]
The producer \path{tangent_block_transport.py} checks every equality,
the internal inverse words and the original-seed comparison.

\emph{Two rank-two kernels from one tangency.}\par
We use an elementary finite-index propagation argument. In
\(F=\langle r,u\rangle\), put \(h=(ur)^2\) and
\(v=r^{-1}ur\). Suppose a subgroup \(J\) is normalized by \(h\)
and contains \(u,v\). The identity \(h=ur^2v\) shows that
\(\langle h,u,v\rangle=\langle r^2,u,v\rangle\), the index-two
subgroup of \(F\) consisting of words with even \(r\)-exponent.
This group normalizes \(J\). The normal closure of \(u,v\) in it
is exactly \(\ker(F\to\langle r\rangle)\), as its displayed
Schreier basis proves. Thus that entire kernel is contained in \(J\).

The actual certificate \path{tangent_node_strips.py} reevaluates all
five hundred eighty stored finite orbit words and supplies explicit
expansions for
\(B^{-1}U_2B,B^{-1}U_6B\in I\). It also verifies literally
\[
 \alpha_5|_{\langle B,U_6\rangle}
       =\operatorname{Inn}((U_6B)^2),\qquad
 \alpha_0(B)=A,\quad\alpha_0(U_6)=U_6.
 \tag{B.81.12}
\]
Apply the preceding argument to \(I\cap\langle B,U_6\rangle\).
The upper bound from the filling and transport by \(\alpha_0\) give
\[
\begin{split}
 I\cap\langle B,U_6\rangle
       &=\ker(\langle B,U_6\rangle\to\langle B\rangle),\\
 I\cap\langle A,U_6\rangle
       &=\ker(\langle A,U_6\rangle\to\langle A\rangle).
\end{split}
 \tag{B.81.13}
\]
This uses a whole invariant rank-two subgroup, rather than extrapolating
from finitely many meridians.

\emph{A second whole rank-five kernel from the four nodes.}\par
Put
\[
 \begin{gathered}
 V_1=U_6^{-1}U_1U_6,\quad V_4=U_6^{-1}U_4U_6,\quad V_6=B^{-1}U_6B,\\
 V_2=B^{-1}U_6^{-1}U_2U_6B,\qquad
 V=\langle V_1,V_2,V_4,V_6\rangle.
 \end{gathered}
 \tag{B.81.14}
\]
All four words belong to \(I\): the last two use the just certified
adjacent conjugates. Together with \(A,B,C\) they form a free basis;
both inverse maps are included in the same exact certificate.
The four actual node restrictions are
\[
 \alpha_k|_{\langle C,V_j\rangle}
       =\operatorname{Inn}(V_jC),\qquad
        (k,j)=(2,1),(3,6),(4,4),(5,2).
 \tag{B.81.15}
\]
For each pair let \(J_j=I\cap\langle C,V_j\rangle\).
It is normalized by \(V_jC\) and contains \(V_j\), so it is
normalized by \(C\) as well. Therefore it contains every
\(C^nV_jC^{-n}\). Combining the four indices and using the upper
bound \(I\leq R\) proves
\[
 I\cap\bigl(V*\langle C\rangle\bigr)
       =\ker\bigl(V*\langle C\rangle\longrightarrow\langle C\rangle\bigr).
 \tag{B.81.16}
\]
This is another exact rank-five free-factor intersection. The independent
checker \path{work_CE/code/verify_tangent_strips.py}, with its genuine
output and \path{work_CE/tangent_strips_independent_review.json}, imports
neither the producing programs nor their folding code. It independently
reevaluates all five hundred eighty source words, both basis changes,
the stored adjacent-conjugate expansions and every action restriction
used in (B.81.10)--(B.81.16). All checks pass.

\emph{Actual full completions and the remaining gate.}\par
The bundle kernel \(F_7\) carries its full induced topology in
\(F_7\rtimes F_6\): characteristic finite-index subgroups are
cofinal in the finitely generated fiber, and the action on each
finite characteristic quotient has finite image. These quotients,
together with finite quotients of the free base, give the exact
completed semidirect product. Similarly the source free retraction
identifies its actual completed kernel with the closure of its
discrete kernel. Its compact image is therefore \(\overline I\).

Every free factor in (B.81.11), (B.81.13) and (B.81.16) has full
induced topology, by its explicit free retraction. Its discrete
filling kernel is dense in the corresponding closed completed
kernel, as follows from the same quotient presentation. The lower
inclusions just proved and the upper bound
\(\overline I\leq\ker(\widehat F_7\to\widehat F_3)\) therefore
prove exactly the same intersection formulas with hats and actual
closures. No abstract completion of any infinite-rank kernel is
identified with an ambient subgroup.

These formulas do not show that \(A,B,C\), or even \(D\), normalize
all of \(I\). The literal infinity calculation illustrates the
remaining issue: its retained quotient agrees with the expected
composition, but its actual image of \(D\) has cyclic core
\(CkAk^{-1}\), where
\(k=U_6^{-1}U_4^{-1}U_6U_1^{-1}U_2U_6\), rather than the length-two
core of \(AC\). The actual correction words are retained in
\path{tangent_infinity_action.json}. Uniform removal of such
corrections on all source translates, the full equality \(I=R\)
or \(\overline I=\ker(\widehat F_7\to\widehat F_3)\), and the
later joint-source comparison remain unresolved.
\end{proof}
